 \begin{center}
        {\large \textbf{Вариант №22 \\
        }}
    \end{center}

    \textbf{\underline{Цель работы:} Изучить основы работы с процессами в ОС Linux, а также освоить сигналы как универсальное средство межпроцессного взаимодействия.
    } 

    \begin{enumerate}
        \item \hyperref[sec:1]{Изучить основные понятия о процессах в ОС Linux, относящиеся к их созданию и управлению.}
        
        \item \hyperref[sec:2]{Ознакомиться с базовыми системными вызовами POSIX для работы с процессами: \texttt{fork()}, \texttt{wait()}, \texttt{waitpid()}.}
        
        \item \hyperref[sec:3]{Изучить назначение и возможности сигналов в Linux, их классификацию (системные, пользовательские) и обработку.}
        
        \item \hyperref[sec:4]{Освоить установку обработчиков сигналов с помощью \texttt{signal()} и \texttt{sigaction()}, а также способы отправки сигналов с помощью \texttt{kill()} и \texttt{sigqueue()}.}
        
        \item \hyperref[sec:5]{Познакомиться с соответствующими консольными инструментами для мониторинга и управления процессами и сигналами.}
        \begin{enumerate}[label=\arabic{enumi}.\arabic*]
            \item \hyperref[sec:5.1]{Изучение команд \texttt{ps}, \texttt{top}, \texttt{htop} — просмотр списка процессов, их состояний и ресурсов.}
            \item \hyperref[sec:5.2]{Изучение команд \texttt{jobs}, \texttt{bg}, \texttt{fg} — управление фоновыми и приостановленными задачами в оболочке.}
            \item \hyperref[sec:5.3]{Изучение команд \texttt{kill}, \texttt{pkill}, \texttt{killall} — отправка сигналов процессам по PID или имени.}
            \item \hyperref[sec:5.4]{Изучение команды \texttt{ulimit} — просмотр и установка ресурсных ограничений (например, число процессов, размер стека).}
            \item \hyperref[sec:5.5]{Изучение команды \texttt{strace} — отслеживание системных вызовов процесса, включая работу с сигналами.}
        \end{enumerate}
        
        \item \hyperref[sec:6]{Выполнить индивидуальное задание для закрепления знаний на практике, подготовив соответствующую программу на языке С (номер задания соответствует номеру студента по журналу; если этот номер больше, чем максимальное число заданий, тогда вариант задания вычисляется по формуле: номер по журналу \% максимальный номер задания, где \% — остаток от деления; если остаток равен 0, студенту назначается задание с максимальным номером).}
        
        \item \hyperref[sec:7]{Подготовить отчёт по работе, который должен включать: краткое описание всех использованных системных вызовов и команд; текст программы; протоколы; скриншоты с демонстрацией поведения ОС Linux; выводы.}
    \end{enumerate}